\section{Main}\label{sec:introduction}

% Opening: The assumption and the problem
The past decade of genome-wide association studies (GWAS) has operated under an implicit assumption: larger sample sizes will yield proportionally more biological insight. Biobank-scale efforts have delivered on this promise for many traits, with the number of trait-associated loci growing steadily with statistical power \cite{abdellaoui_15_2023}. Yet new loci do not automatically translate into new biological understanding. A locus may implicate a gene already known to be relevant through other associations, or it may refine the evidence for a mechanism already captured by existing functional annotations. The field has lacked a systematic framework for distinguishing genuine biological discovery from incremental refinement of existing knowledge---and therefore for diagnosing when GWAS investment has reached the point of diminishing returns.

% The three levels of discovery
This gap arises because ``genetic discovery'' conflates at least three distinct processes: \textit{variant discovery} (identifying new associated loci), \textit{gene discovery} (identifying new disease-relevant genes), and \textit{mechanism discovery} (identifying new biological pathways or processes). These processes need not saturate in parallel. A trait may continue to yield new loci while the genes they implicate are already functionally connected to known associations. Conversely, a trait may have few loci but each representing a genuinely independent biological mechanism. Disentangling these trajectories requires separating what GWAS directly detects---statistical associations---from what those associations reveal about gene-trait relevance when interpreted through functional annotations.

% Brief PIGEAN reference and what we do
We recently developed PIGEAN \cite{}, a Bayesian framework that separately quantifies \textit{direct genetic support} (the probability that a gene carries trait-associated variants) and \textit{indirect genetic support} (the probability that a gene is trait-relevant based on shared functional annotations with associated genes). Here, we leverage this decomposition to trace the accumulation of gene discovery as a function of GWAS power across {{NUM_TRAITS_SATURATION_ANALYSIS}} common traits, revealing that a substantial fraction of the phenome has reached or is approaching saturation for gene-level discovery---even as variant-level discovery continues.
